\documentclass{article}
\usepackage{amssymb, amsmath}
\usepackage[left=1.5cm, right=2cm]{geometry}
\begin{document}

\textbf{Public Key Cryptography}

\textbf{Public Key Cryptography}

Kerckhoff's Principle: the cryptosystem (algorithm) is public\\
\textit{Symmetric or shared key}: only the key is secret, used both for encryption and decryption. Shared keys for MACs, PRFs, etc..\\
\textit{Asymmetric} cryptographic capability: encryption capability without decryption capability.

\textbf{Public Key Cryptosystem (PKC}

Kerckhoff: cryptosystem (algorithm) is public\\
DH76: can encryption key be public too?\begin{itemize}
	\item Decryption key will be different (and private)
	\item Everybody can send me emails, but only I can read them.
\end{itemize}

\textbf{Is it Only About Encryption?}

Also: digital signatures for integrity and non-repudiation.\\
Sign with private key $s$, verify with public key $v$.\\
Recall MACs: a shared key cryptosystem for message authentication\\
Private signing key, public validation key

\textbf{More: Key-Exchange Protocols}

Establish shared key between Alice and Bob without assuming an existing shared (master) key\\
Use public information from Alice and Bob to setup shared key $k$\\
Eavesdroppers cannot learn the key $k$

\textbf{Public keys solve more problems}

Signatures provide evidence (everyone can validate, only owner can sign)\\
Established shared secret keys\begin{itemize}
	\item Use authenticated public keys, signed by trusted certificate authority
	\item Or use DH (Diffie Hellman) key exchange
\end{itemize}
Stronger resiliency to key exposure\begin{itemize}
	\item Perfect forward secrecy and recover security
	\item Threshold security: resilient to key exposure of $t$ out of $n$ parties
\end{itemize}

\textbf{Public keys are easier...}
To distribute: from directory or from incoming message (still require authentication), and less keys to distribute\\
To maintain: can keep in non-secure storage as long as being validated before using. Also less keys: $O(|parties|)$ versus $O(|parties|^2)$\\
So why not always use public key crypto?

\textbf{The Price of PKC}

Assumptions:
Applied PKC algorithms are based on a small number of specific computational assumptions\begin{itemize}
	\item Mainly: hardness of factoring and discrete-log
	\item Both may fail against quantum computers
\end{itemize}
Overhead:\begin{itemize}
	\item Computational
	\item Key length (public key is long)
	\item Output length (e.g. ciphertext or signature)
\end{itemize}

\textbf{Public key crypto is harder...}

Requires related public, private keys\begin{itemize}
	\item Usually we say a keypair (pk, sk)
	\item Public key does not expose private key
\end{itemize}
Substantial overhead\begin{itemize}
	\item Successful cryptanalytic shortcuts $\rightarrow$ need long keys
	\item Elliptic curves (EC) may allow shorter keys (almost no shortcuts found)
	\item Complex computations, e.g. complex (slow) keygen
\end{itemize}

\textbf{In summary}

Minimize use of PKC\\
If possible, minimize length of the input\\
How?\\
For signatures, hash-then-sign\\
For public-key encryption, hybrid encryption

\textbf{Hybrid Encryption}

Challenge: public key cryptosystems are slow\\
Hybrid encryption:\begin{itemize}
	\item Use a shared key encryption scheme to encrypt all messages,
	\item But use a public key encryption scheme to exchange the shared key\begin{itemize}
			\item Alice generates $k$, encrypts it under Bob's public key and sends the ciphertext $c_k$ to Bob.
			\item Bob can decrypt and recover $k$, and use $k$ to decrypt $c_M$
		\end{itemize}
\end{itemize}

\textbf{Number Theory Review}

$(a+b) \mod m = [(a \mod m) + (b \mod m)] \mod m$

\textbf{Multiplicative Inverse}

To compute $a/c \mod m$, multiply $a$ by the multiplicative inverse of $c$\\
That is, compute $a/c\mod m = ac^{-1}\mod m$\\
Where $c^{-1}$ is the multiplicative inverse such that $cc^{-1}\mod m=1$\\
Not all ints have multiplicative inverses with respect to a specific modulus $m$.

\textbf{Modular Exponentiation}

Encountered a lot: discrete log scheme, RSA, etc\\
Fermat's Little Theorem (for prime modulus):
\[
	a^b \mod p = (a\mod p)^{b\mod(p-1)} \mod p
\]

\textbf{Group Theory Review}

A group G is a cyclic if there is a an element $g \epsilon G$ s.t. for every element $a\epsilon G$, there is an int $i$ s.t. $a=g^i$. Such an element $g$ is called a generator of $G$. The order of $G$ is the total.

\textbf{Key Exchange}

\textbf{Defining a Key Exchange Protocol}

Must satisfy:\\
Correctness: both parties compute the same shared key\\
Key Indistinguishability: the key that the two parties establish is indistinguishable from random

\textbf{Discrete Log Problem}

A computationally hard problem that is hard to solve but easy to verify\\
Given a generator $g$ and an element $a\epsilon G$, find $i$ s.t. $a=g^i$. Verification: exponentiation (efficient algorithm)\\
Computing log is quite efficient over the real numbers. But is discrete-log hard?\begin{itemize}
	\item Weak, where it's not hard: prime $p$ where $(p-1)$ has only 'small' prime factors, mistakes/trapdoors found
	\item Safe-prime groups: $p=2q+1$, where $q$ is sufficiently large
\end{itemize}
Group needs to be big, cyclic with generator, and discrete log problem within group to be hard.\\
Game: pick random element in group. First, pick the power (random int) and then compute the element as the generator to that power. Attacker is given $g^y$, where they must compute $y$.\\
$y$ is the secret key, the public key is $g^y$.

\textbf{The Diffie-Hellman (DH) Key Exchange Protocol and The Computational/Decisional Diffie-Hellman Assumptions (CDH/DDH)}

\textbf{DH Key-Exchange}

Setup: agree on a random safe (large) prime $p$ and generator $g$ for the cyclic multiplicative group $Z_p^*$.\\
Alice: pick at random a secret int $a$ from $Z_p^*$, then compute $P_A=g^a\mod p$, and send $P_A$ to Bob.\\
Bob: pick at random secret int $b$ from $Z_p^*$, then compute $P_B=g^b\mod p$, and send $P_B$ to Alice.\\
Both parties: compute the shared key $k=g^{ab}\mod p$\\
$(g^b)^a=g^{ab}=(g^a)^b\mod p$; Alice raises Bob's $g^b$ to the $a^{th}$ power, and vice versa.\\
Correctness: two parties establish identical, long-enough, and random key.

\textbf{Caution: Auth the Public Keys!}

DH Key-Exchange is only secure against eavesdroppers, but not MitM attackers (since they can change the messages)\\
So, the public messages being sent must be authenticated, e.g. using digital signatures.

\textbf{Security of DH Key Exchange}

Assume authenticated comms.\\
DH key exchange requires stronger assumption than discrete log\\
The Computational Diffie-Hellman (CDH) Assumption: given $g^b\mod p$ and $g^a\mod p$, an efficient adversary cannot compute $g^{ab}\mod p$ with non-negligible probability.\\
So DH key exchange protocol is secure for groups in which the CDH assumptions holds.

\textbf{Using DH Securely?}

Can $g^a,\ g^b$ expose \textit{something} about $g^{ab}\mod p$? Finding at least on ebit about it is easy\\
Option 1: use DH but with a stronger group where $g^a,g^b,g^c$ are indistinguishable\\
Option 2: use DH with Z and safe prime (where only CDH holds), but use a key derivation function to derive a secure shared key (like a randomness-extracting hash function).

\end{document}
